\chapter{Conclusions}\label{ch:conclusions}
A reduced temporal model evolves O, ozone, H, OH, HO$_2$, H$_2$O$_2$ and $\Delta$ at 51 levels over 50--100~km: 357 ODEs. O($^1$D), $B_0$ and $B_1$ remain algebraic; $R_H=\mathrm{OH}+\mathrm{HO_2}$ is diagnostic. Accepted event rates and stoichiometry are retained.

The temporal extension demonstrates that the original stationary radical/peroxide reduction is not globally compatible with the tested illumination history. The OH partition loses a physical root at dawn, and peroxide elimination becomes singular in a dark limit with positive production and zero algebraic loss. Promoting those species dynamically resolves these reduction-domain failures without changing reactions, forcing or constants.

UTC date, latitude and longitude determine solar geometry. Spherical shells determine shadow and attenuation; ultraviolet opacity responds to evolving O/O$_3$. Complete selected A0/B/IRA line sets and historical CIA are retained. The tested advanced A-band sensitivity is small downstream.

The finite-horizon model is numerically resolved under the tested controls. Maximum relevant all-species differences are approximately $0.424\%$ for base/tighter BDF and $0.00946\%$ for tighter BDF/Radau. Independent background and NIR interpolation refinements also meet the practical approximately $0.5\%$ output target. Physical states, remaining QSSA, event budgets and exact solar shadow pass the accepted checks.

The reference dawn shows upper-level ozone decreases, singlet growth and non-monotonic ozone near 80~km. Season/latitude responses combine geometry, background and initialization effects. With identical initial states, frozen/dynamic atmosphere differences reach $69.21\%$ in ozone and $24.61\%$ in singlet oxygen.

Initialization remains a material physical input: atomic-state scaling produces differences of tens of percent despite consistent fast-state initialization. Date/location cannot determine a validated chemical state; the reference-noon route remains approximate and explicit-state input necessary.

The accepted forward, finite-horizon objective is achieved. Observational validation and satellite retrieval remain future work, requiring improved initial profiles, transport assessment and compatible measurements.
